Brazil's Unified Health System (SUS) gives every resident the right to public hospital care.
In practice, hospitals and medical specialties are spread unevenly across the country, so many patients leave their home town to be admitted elsewhere.
Knowing where these patients go, and for which specialties, helps managers plan hospital beds and find regions that depend on care produced somewhere else.

The Hospital Information System of SUS (SIH/SUS), published by DATASUS, records every publicly funded hospital admission.
Each record stores the patient's municipality of residence and the municipality of the hospital, so a single admission already describes a movement between two places.
The data is large: our extract holds more than 180 million admissions from 2008 to 2023.
At this size, tables and static maps make it hard to see the overall pattern of movement between cities and states.

We present an interactive flow map of these hospitalization flows.
Each arc links the patient's residence to the hospital municipality, and its width grows with the number of admissions.
The map works at three levels, from the whole country down to one state and then one city, so the user opens detail only where it is needed.
Filters by medical specialty, year range, and month narrow the data, and a side panel ranks the cities that most attract patients and the ones whose residents most depend on care elsewhere.

The design answers four questions raised in the first part of this work:
(1) what share of patients is treated outside their home municipality, and how it varies by specialty;
(2) where the care gaps are for specific specialties such as oncology;
(3) which cities sustain the hospital care of a whole region;
and (4) whether seasonal peaks in admissions are driven by patients coming from outside.
This paper describes how we prepared the data and designed the visualization, and shows how it answers these questions for admissions between 2008 and 2023.
